\subsection{The fundamental theorem of Galois theory}

THEOREM 4. --- Let N be a Galois extension of K and \(\Gamma\) its Galois group. Let \(\mathcal{K}\) be the set of subextensions of N and \(\mathcal{G}\) the set of closed subgroups of \(\Gamma\) . For every subgroup \(\Delta \in \mathcal{G}\) we denote by \(k(\Delta)\) the field of invariants of \(\Delta\) and for every subfield E \(\in \mathcal{K}\) we denote by g(E) the group of E-automorphisms of N. Then \(\Delta \mapsto k(\Delta)\) is a bijection of \(\mathcal{G}\) onto \(\mathcal{K}\) , and E \(\mapsto\) g(E) is the inverse bijection.

\begin{enumerate}
\def\labelenumi{\Alph{enumi})}
\tightlist
\item
  The relation \(E = k(g(E))\) (for \(E \in \mathcal{K}\) ) is a consequence of the following more precise lemma:
\end{enumerate}

Lemma 1. --- Let E be a subextension of N. Then N is a Galois extension of E and Gal(N/E) is a closed subgroup of Gal(N/K) with the induced topology.

Let \(x \in \mathbb{N}\) ; the minimal polynomial \(f\) of \(x\) over \(\mathbf{E}\) divides in \(\mathbf{E}[\mathbf{X}]\) the minimal polynomial \(g\) of \(x\) over \(\mathbf{K}\) (V, p.~17, Cor. 2). Since \(\mathbf{N}\) is Galois over \(\mathbf{K}\) , the polynomial \(g\) is a product in \(\mathbf{N}[\mathbf{X}]\) of distinct factors of degree 1; hence the same is true of \(f\) and so \(\mathbf{N}\) is Galois over \(\mathbf{E}\) .

Let \(\Gamma\) be the Galois group of N over K and \(\Delta\) that of N over E. By definition \(\Delta\) is the subgroup of \(\Gamma\) consisting of all \(\sigma\) such that \(\sigma(x) = x\) for all \(x \in E\) . Now for each \(x \in E\) the mapping \(\sigma \mapsto \sigma(x)\) of \(\Gamma\) into the discrete space N is continuous, hence \(\Delta\) is closed in \(\Gamma\) . Let \(\sigma \in \Gamma\) ; for \(x_1, \ldots, x_n\) in N let U \((x_1, \ldots, x_n)\) be the set of all \(\tau \in \Gamma\) such that \(\tau(x_i) = \sigma(x_i)\) for \(1 \leqslant i \leqslant n\) ; put

\[
\mathrm{V} \left(x _ {1}, \dots , x _ {n}\right) = \mathrm{U} \left(x _ {1}, \dots , x _ {n}\right) \cap \Delta .
\]

Then the family of sets \(\mathrm{U}(x_{1},\ldots,x_{n})\) (resp. \(\mathrm{V}(x_{1},\ldots,x_{n})\) ) is a base of neighbourhoods of \(\sigma\) in \(\Gamma\) (resp. \(\Delta\) ). Hence the topology on \(\Delta\) is that induced by \(\Gamma\) .

\begin{enumerate}
\def\labelenumi{\Alph{enumi})}
\setcounter{enumi}{1}
\tightlist
\item
  The relation \(\Delta = g(k(\Delta))\) (for \(\Delta \in \mathcal{G}\) ) is a consequence of the following more precise lemma:
\end{enumerate}

Lemma 2. --- Let \(\Delta\) be a subgroup of \(\Gamma\) . Let \(E\) be the field of invariants of \(\Delta\) ; then the Galois group of \(N\) over \(E\) is the closure of \(\Delta\) in \(\Gamma\) .

The Galois group of N over E is closed in \(\Gamma\) (Lemma 1) and contains \(\Delta\) , hence it contains the closure \(\bar{\Delta}\) of \(\Delta\) . Let \(\sigma\) be an E-automorphism of N and let \(x_{1}, \ldots, x_{n}\) be in N. Since N is Galois over E (Lemma 1) there exist (V, p.~57, Prop. 2) a subextension \(N_{0}\) of N, Galois of finite degree over E and containing \(x_{1}, \ldots, x_{n}\) . Let \(\Delta_{0}\) be the image of the subgroup \(\Delta\) of \(\operatorname{Gal}(N/E)\) under the restriction homomorphism of \(\operatorname{Gal}(N/E)\) into \(\operatorname{Gal}(N_{0}/E)\) . Since \([N_{0}:E]\) is finite, Dedekind's theorem (V, p.~27, Cor. 2) shows that \(\operatorname{Gal}(N_{0}/E)\) is finite. Hence \(\Delta_{0}\) is finite, and since E is the field of invariants of \(\Delta_{0}\) , we have \(\Delta_{0}=\)

Gal(N \(_0\) /E) (V, p.~66, Th. 3). In particular, \(\Delta_0\) contains the restriction of \(\sigma\) to N \(_0\) . Therefore there exists \(\tau \in \Delta\) such that \(\sigma\) and \(\tau\) have the same restriction to N \(_0\) , whence \(\sigma(x_1) = \tau(x_1)\) , \ldots, \(\sigma(x_n) = \tau(x_n)\) . It follows that \(\sigma\) is a limit point of \(\Delta\) in \(\Gamma\) , and hence Gal(N/E) \(\subset \bar{\Delta}\) .

COROLLARY 1. --- Let E and E' be two subfields of N containing K; then E ⊂ E' if and only if g(E) ⊃ g(E'). If Δ and Δ' are two closed subgroups of Γ, then Δ ⊂ Δ' if and only if k(Δ) ⊃ k(Δ').

For the two inverse bijections \(E \mapsto g(E)\) and \(\Delta \mapsto k(\Delta)\) are inclusion-reversing.

COROLLARY 2. --- Let \((\mathrm{E}_i)_{i\in I}\) be a family of subfields of N containing K; put \(\mathbf{L} = \bigcap_{i\in I}\mathbf{E}_i\) and \(\mathbf{M} = \mathbf{K}\left(\bigcup_{i\in I}\mathbf{E}_i\right)\) . Then \(g(L)\) is the smallest closed subgroup of \(\Gamma\)

containing \(\bigcup_{i\in I}g(E_i)\) and we have \(g(M) = \bigcap_{i\in I}g(E_i)\) .

The first assertion follows from Cor. 1 and the second is immediate.

COROLLARY 3. --- For \(i=1\) , 2 let \(\mathrm{E}_{i}\) be a subfield of N containing K and let \(\Delta_{i}=g(\mathrm{E}_{i})\) . For any \(\sigma\in\Gamma\) the relations \(\sigma(\mathrm{E}_{1})=\mathrm{E}_{2}\) and \(\sigma\Delta_{1}\sigma^{-1}=\Delta_{2}\) are equivalent.

For we have \(\tau \in g(\sigma(E_1))\) if and only if \(\tau\sigma(x) = \sigma(x)\) , that is, \(\sigma^{-1}\tau\sigma(x) = x\) , for all \(x \in E_1\) ; this amounts to saying that \(\sigma^{-1}\tau\sigma \in \Delta_1\) , whence \(g(\sigma(E_1)) = \sigma\Delta_1\sigma^{-1}\) .

COROLLARY 4. --- Let E be a subfield of N containing K and let \(\Delta = g(E)\) . For E to be Galois over K it is necessary and sufficient that \(\Delta\) should be a normal subgroup of \(\Gamma\) . When this is so, the restriction homomorphism of \(\Gamma\) into Gal(E/K) defines by passage to quotients a topological group isomorphism of \(\Gamma/\Delta\) onto Gal(E/K).

Since N is separable over K, the same is true of E (V, p.~36, Prop. 1). Therefore E is Galois over K if and only if it is quasi-Galois over K; this also means that \(\sigma(E) = E\) for every K-automorphism \(\sigma\) of N (V, p.~52, Prop. 1 and p.~54, Prop. 3). By Cor. 3 this is equivalent to \(\sigma\Delta\sigma^{-1} = \Delta\) for all \(\sigma \in \Gamma\) .

The restriction homomorphism \(\varphi: \operatorname{Gal}(N/K) \to \operatorname{Gal}(E/K)\) is continuous and surjective (V, p.~60, Prop. 3) and its kernel is clearly equal to \(\Delta = \operatorname{Gal}(N/E)\) . Since \(\Gamma\) is compact, the homomorphism of \(\Gamma/\Delta\) onto \(\operatorname{Gal}(E/K)\) derived from \(\varphi\) by passage to quotients is an isomorphism of topological groups (Gen.~Top., I, p.~87, Cor. 2).

COROLLARY 5. --- Let E be a subfield of N containing K. For E to have finite degree over K it is necessary and sufficient that \(g(E)\) should be open in \(\Gamma\) . When this is so, the index \((\Gamma:g(E))\) is finite and equal to \([E:K]\) .

For \(g(E)\) to be open it is necessary and sufficient that there should exist a subextension F of N, of finite degree over K, such that in the notation of V, p.~60, \(g(E)\) contains \(U_{F}(Id_{N}) = g(F)\) . The relation \(g(E) \supset g(F)\) is equivalent to \(E \subset F\) by Cor. 1 (V, p.~68), whence the first assertion of Cor. 5.

Suppose that \([E:K]\) is finite. Let \(\Omega\) be an algebraic closure of K containing N as subextension (V, p.~23, Th. 2) and let H be the set of K-homomorphisms of E into \(\Omega\) . Every element of H is induced by a K-automorphism of \(\Omega\) (V, p.~52, Prop. 1), and since N is quasi-Galois over K, the mapping \(\sigma\mapsto\sigma|E\) of \(\Gamma\) into H is surjective. For \(\sigma\) and \(\sigma'\) in \(\Gamma\) to have the same restriction to E it is necessary and sufficient that \(\sigma^{-1}\sigma'\in g(E)\) , whence Card \(H=(\Gamma:g(E))\) . Finally since E is an etale algebra over K, we have Card \(H=[E:K]\) (V, p.~32, Prop. 4), so in conclusion we have \((\Gamma:g(E))=[E:K]\) .

COROLLARY 6. --- For \(i=1,2\) let \(\mathbf{E}_{i}\) be a subextension of N and \(\Gamma_{i}\) the Galois group of N over \(\mathbf{E}_{i}\) . The following conditions are equivalent:

\begin{enumerate}
\def\labelenumi{\alph{enumi})}
\item
  The group \(\Gamma\) is the direct product of the subgroups \(\Gamma_1\) and \(\Gamma_2\) .
\item
  The extensions \(\mathbf{E}_1\) and \(\mathbf{E}_2\) are Galois over \(\mathbf{K}\) , we have \(\mathbf{E}_1 \cap \mathbf{E}_2 = \mathbf{K}\) and
\end{enumerate}

\[
\mathbf {K} \left(\mathrm{E} _ {1} \cup \mathrm{E} _ {2}\right) = \mathbf {N}.
\]

For \(\Gamma\) to be the direct product of the subgroups \(\Gamma_{1}\) and \(\Gamma_{2}\) it is necessary and sufficient that the following conditions hold (I, p.~48, Prop. 15):

\begin{enumerate}
\def\labelenumi{(\roman{enumi})}
\item
  the subgroups \(\Gamma_{1}\) and \(\Gamma_{2}\) are normal in \(\Gamma\) ;
\item
  \(\Gamma_1 \cap \Gamma_2 = \{\varepsilon\}\) , where \(\varepsilon\) is the neutral element of \(\Gamma\) ;
\item
  \(\Gamma = \Gamma_{1} \cdot \Gamma_{2}\) .
\end{enumerate}

Now (i) means that \(E_{1}\) and \(E_{2}\) are Galois over K (Cor. 4). By Cor. 2, condition (ii) is equivalent to \(\mathbf{N} = \mathbf{K}(\mathbf{E}_{1} \cup \mathbf{E}_{2})\) , and finally if (i) and (ii) hold, \(\Gamma_{1}\Gamma_{2}\) is the least subgroup of \(\Gamma\) containing \(\Gamma_{1} \cup \Gamma_{2}\) ; it is closed because \(\Gamma_{1}\) and \(\Gamma_{2}\) are compact and the mapping \((\sigma, \tau) \mapsto \sigma\tau\) of \(\Gamma_{1} \times \Gamma_{2}\) into \(\Gamma\) is continuous (Gen.~Top., I, p.~63, Cor. 1). Cor. 2 now shows (iii) to be equivalent to \(E_{1} \cap E_{2} = K\) , and this proves the equivalence of a) and b).

Remark. --- With the notation of Cor. 6 suppose that conditions a) and b) hold. The restriction homomorphisms \(\varphi_{i}:\Gamma\to\mathrm{Gal}(\mathrm{E}_{i}/\mathrm{K})\) for \(i=1,2\) induce topological group isomorphisms

\[
\Psi_ {1}: \Gamma_ {2} \rightarrow \operatorname{Gal} \left(\mathrm{E} _ {1} / \mathrm{K}\right), \quad \Psi_ {2}: \Gamma_ {1} \rightarrow \operatorname{Gal} \left(\mathrm{E} _ {2} / \mathrm{K}\right).
\]

By \(a\) we see that the mapping \(\sigma \mapsto (\varphi_1(\sigma), \varphi_2(\sigma))\) is a topological group isomorphism of \(\operatorname{Gal}(\mathbf{N}/\mathbf{K})\) onto \(\operatorname{Gal}(\mathrm{E}_1/\mathrm{K}) \times \operatorname{Gal}(\mathrm{E}_2/\mathrm{K})\) .
% semantic-fragments: legacy-input-seam
\relax{}
